% ECONOMICS_PROBLEM: {"id": "D63-177", "title": "Polynomial-time computation of simultaneous EEFX and EFL allocations", "jel_code": "D63", "source_id": "OP-0192"}
% !TeX program = xelatex
\documentclass[11pt,a4paper]{article}
\usepackage[margin=23mm]{geometry}
\usepackage{fontspec}
\setmainfont{Latin Modern Roman}
\usepackage{amsmath,amssymb,mathtools}
\usepackage{xcolor,enumitem,booktabs,longtable,array,xurl,hyperref}
\definecolor{muted}{HTML}{53636C}
\hypersetup{colorlinks=true,urlcolor=blue,linkcolor=blue}
\setlength{\parindent}{0pt}
\setlength{\parskip}{5pt}
\newcommand{\R}{\mathbb R}
\newcommand{\E}{\mathbb E}
\newcommand{\N}{\mathbb N}
\newcommand{\ind}{\mathbf 1}
\newcommand{\Var}{\operatorname{Var}}
\newcommand{\Cov}{\operatorname{Cov}}
\newcommand{\supp}{\operatorname{supp}}
\DeclareMathOperator*{\argmin}{arg\,min}
\DeclareMathOperator*{\argmax}{arg\,max}
\newcommand{\entrymeta}[3]{{\sffamily\small\color{muted}\textbf{\texttt{#1}}\quad|\quad #2\par #3\par}}
\newcommand{\Problemref}[1]{\texttt{#1}}
\newcommand{\JELref}[2]{\texttt{#2} (JEL \texttt{#1})}
\begin{document}
\section*{Polynomial-time computation of simultaneous EEFX and EFL allocations}
\textbf{Problem ID:} \texttt{D63-177}\quad \textbf{JEL:} \texttt{D63}\par
% BEGIN ECONOMICS STATEMENT
\label{problem:OP-0192}
% GAME_THEORY_PROBLEM: {"local_id": "GT-062", "original_number": 62, "kind": "P", "entry_kind": "statement_only", "published": false, "review_required": true, "category": "game_theory"}
\label{game:GT-062}
{\sffamily\small\color{muted}
\noindent\textbf{ID: GT-062.}\quad\textbf{Category:} Fair division algorithms.
\quad\textbf{Kind: P.}\quad\textbf{Status:} Open in the cited source.
\par}
\subsection*{Definitions and mathematical statement}
Let \(N=\{1,\ldots,n\}\), \(n\geq 2\), let \(M\) be a finite set of goods,
and let \(v_i(S)=\sum_{g\in S}a_{ig}\) with
\(a_{ig}\in\mathbb Q_{\geq 0}\). A complete allocation is an ordered
partition \(A=(A_i)_{i\in N}\) of \(M\).

Use the following \emph{all-good} convention for epistemic envy-freeness
up to any good (\(\mathrm{EEFX}\)). An allocation is \(\mathrm{EEFX}\)
if, for every \(i\), there is an ordered partition
\((P^{(i)}_j)_{j\in N\setminus\{i\}}\) of \(M\setminus A_i\) such that
\[
 v_i(A_i)\geq v_i(P^{(i)}_j\setminus\{g\})
 \quad\text{for all }j\ne i,\ g\in P^{(i)}_j.
\]
Certificates for different agents may be different.
An allocation is envy-free up to one less-preferred good
(\(\mathrm{EFL}\)) if, for every \(i,j\), either
\[
 \bigl|\{g\in A_j:a_{ig}>0\}\bigr|\leq 1,
\]
or there is \(g\in A_j\) such that
\[
 v_i(A_i)\geq v_i(A_j\setminus\{g\})
 \quad\text{and}\quad
 v_i(A_i)\geq a_{ig}.
\]

\paragraph{Question.}
Is there a deterministic algorithm that, for every such input, outputs
a complete allocation satisfying both properties in time polynomial in
the binary encoding length of \((n,M,(a_{ig}))\)?

\subsection*{Short English statement}
Can an allocation with both EEFX and EFL be found efficiently?

\subsection*{Variants and scope boundaries}
Existence and polynomial-time computation are separate assertions.
The first display tests removal of every good, including zero-value goods.
Entry~\ref{game:GT-063} uses the positive-value convention of a different paper.
Replacing EFL by EF1 changes the computational question.

\subsection*{Source relationship and status}
[S1, Section 4] establishes existence; Section 6 leaves efficient
simultaneous computation open. The statement uses the definitions in the
revised version, which added an author and changed the original title.

\subsection*{References}
\begingroup\footnotesize\raggedright
\noindent[S1] Hannaneh Akrami, Uriel Feige, Ryoga Mahara, Kurt Mehlhorn,
and Nidhi Rathi,
\emph{The Power of Share-Based Notions in Proving Envy-Based Fairness
Guarantees}, 2026, arXiv:2602.11732v2.
\url{https://arxiv.org/html/2602.11732v2}.
Locators: Section 2, Definitions 2.1 and 2.5--2.7;
Section 4; Section 6, final future-directions item.
\par\endgroup

\subsection*{Common conventions: Source mathematical conventions}

\textbf{Finite games.} For a finite set $A$, $\Delta(A)$ is its probability
simplex. Players choose independent mixed strategies in Nash equilibrium;
correlation is allowed only when explicitly part of the solution concept.
In a game with payoff functions $u_i$ and strategy profile $x$, an additive
$\varepsilon$ Nash equilibrium satisfies
\[
 u_i(x)\geq u_i(x_i',x_{-i})-\varepsilon
 \quad\text{for every player $i$ and every admissible deviation $x_i'$}.
\]
For numerical approximation, payoffs are normalized as locally specified.
An $\varepsilon$ payoff residual is distinct from distance at most
$\varepsilon$ to the set of exact equilibria. Point convergence, convergence
of empirical play, and convergence in distance to an equilibrium set are
also distinct.

\textbf{Computational representations.} Unless stated otherwise, a complexity
question takes rational data with binary encoding; polynomial time is measured
in total bit length. A fixed-accuracy polynomial algorithm is not necessarily
polynomial in $1/\varepsilon$, and neither is necessarily polynomial in
$\log(1/\varepsilon)$. Succinct games and value, demand, or sampling oracles
must declare their interfaces. Exact real solutions require an explicit
representation; an unspecified list of arbitrary real numbers is not a
finite computational output. A claimed algorithmic barrier is meaningful
only for its specified model and reductions.

\textbf{Dynamic games.} Histories, information, observation, and allowable
randomization are part of the model. A stationary strategy depends only on
the current state; a positional strategy in a deterministic graph game
chooses one action at each controlled vertex. Nash equilibrium at an initial
state and equilibrium after every history are different requirements.
An expectation of a pathwise $\liminf$ payoff, a $\liminf$ of expected averages,
a limiting expected average, and a uniform finite-horizon equilibrium are
different criteria. None is silently substituted for another.

\textbf{Allocation and mechanisms.} A complete allocation assigns every item
exactly once unless otherwise stated. Zero-valued goods, negative values,
capacity constraints, feasibility, lotteries, and copies of goods affect
fairness definitions and are specified locally. Dominant-strategy incentive
compatibility, Bayesian incentive compatibility, universal truthfulness,
and truthfulness in expectation impose different inequalities. Whenever
an objective ratio has zero denominator, its convention is stated or those
instances are excluded.

\textbf{Infinite-dimensional models.} Expectations, integrals, stochastic
equations, and optimization objectives require their local regularity,
integrability, filtrations, and admissible domains. A backward--forward
mean-field system is a boundary-value problem; it is not an initial-value
semiflow without an additional construction. Long-time, large-population,
and vanishing-noise limits require an order of limits and a convergence
topology. If a minimum need not be attained, the objective is its infimum
or supremum and the approximation requirement is stated separately.
% END ECONOMICS STATEMENT
\end{document}
